% Fuente única de clase: Recursos/2025-11-18.pdf, páginas 1--9.
En la clase 24 construimos clausuras algebraicas y demostramos su unicidad salvo isomorfismo sobre el campo base. Después estudiamos los campos de descomposición de polinomios, una cota para sus grados y su relación con las extensiones normales.

%%%%%%%%%%%%%%%%%%%%%%%%%%%%%%%%%%%%%%%%%%%%%%%%%%%%%%%%%%%%%%%%%%%%%%%%%%%%%%%
\section{Clausuras algebraicas}
%%%%%%%%%%%%%%%%%%%%%%%%%%%%%%%%%%%%%%%%%%%%%%%%%%%%%%%%%%%%%%%%%%%%%%%%%%%%%%%
%%%%%%%%%%%%%%%%%%%%%%%%%%%%%%%%%%%%%%%%%%%%%%%%%%%%%%%%%%%%%%%%%%%%%%%%%%%%%%%
\subsection{Ausencia de extensiones algebraicas propias}
%%%%%%%%%%%%%%%%%%%%%%%%%%%%%%%%%%%%%%%%%%%%%%%%%%%%%%%%%%%%%%%%%%%%%%%%%%%%%%%
% P01--P02.
Recordemos que un campo $K$ es algebraicamente cerrado si todos los irreducibles de $K[t]$ tienen grado uno; equivalentemente, todo polinomio no constante tiene una raíz en $K$ y se descompone en factores lineales. También equivale a que los ideales maximales de $K[t]$ sean todos de la forma $(t-c)$, con $c\in K$: este anillo es un DIP y sus ideales maximales están generados por irreducibles. Es el mismo argumento aplicado a $\mathbb C$ en el \cref{ejer:C13-complejos}.

\begin{lema}\label{lema:C24-cerrado}
Para un campo $K$ son equivalentes:
\begin{enumerate}[nosep]
\item $K$ es algebraicamente cerrado;
\item toda extensión algebraica de $K$ coincide con $K$;
\item si $K\subseteq L$ es cualquier extensión y $\alpha\in L$ es algebraico sobre $K$, entonces $\alpha\in K$.
\end{enumerate}
\end{lema}
\begin{proof}
Si $K$ es algebraicamente cerrado y $L/K$ es algebraica, el mínimo de cada $\alpha\in L$ es lineal: tiene la forma $t-c$, con $c\in K$. Por tanto $\alpha=c$ y $L=K$.

La segunda condición aplicada a la extensión algebraica $K(\alpha)/K$ da la tercera. Finalmente, sea $f\in K[t]$ irreducible. En el campo $K[t]/(f)$, la clase $\alpha=t+(f)$ es una raíz algebraica de $f$. Por la tercera condición pertenece a $K$, luego $t-\alpha$ divide a $f$ en $K[t]$ y su irreducibilidad obliga a que tenga grado uno.
\end{proof}

\begin{defi}
Una \emph{clausura algebraica} de un campo $k$ es una extensión algebraica $k\subseteq\overline k$ en la que $\overline k$ es algebraicamente cerrado.
\end{defi}
Así, no sólo se descomponen los polinomios de $k[t]$, sino también los que tienen coeficientes en $\overline k$. La condición de algebraicidad impone una minimalidad: si $k\subseteq E\subsetneq\overline k$, entonces $E$ no puede ser algebraicamente cerrado, pues $\overline k/E$ sería una extensión algebraica propia.

\begin{ejer}[Las raíces generan la clausura]\label{ejer:C24-raices-clausura}
Sea $k\subseteq\overline k$ una clausura algebraica. Probar que si un subcampo $E\subseteq\overline k$ contiene a $k$ y a todas las raíces, en $\overline k$, de todos los polinomios no constantes de $k[t]$, entonces $E=\overline k$.
\end{ejer}
\begin{sol}
Todo $\alpha\in\overline k$ es algebraico sobre $k$, por lo que es raíz de su mínimo no constante en $k[t]$. La hipótesis exige que esa raíz pertenezca a $E$. Así $\overline k\subseteq E$, y la otra inclusión ya estaba dada.
\end{sol}

%%%%%%%%%%%%%%%%%%%%%%%%%%%%%%%%%%%%%%%%%%%%%%%%%%%%%%%%%%%%%%%%%%%%%%%%%%%%%%%
\section{Existencia y unicidad de las clausuras}
%%%%%%%%%%%%%%%%%%%%%%%%%%%%%%%%%%%%%%%%%%%%%%%%%%%%%%%%%%%%%%%%%%%%%%%%%%%%%%%
%%%%%%%%%%%%%%%%%%%%%%%%%%%%%%%%%%%%%%%%%%%%%%%%%%%%%%%%%%%%%%%%%%%%%%%%%%%%%%%
\subsection{Adjunción simultánea de raíces}
%%%%%%%%%%%%%%%%%%%%%%%%%%%%%%%%%%%%%%%%%%%%%%%%%%%%%%%%%%%%%%%%%%%%%%%%%%%%%%%
% P02--P04.
\begin{lema}\label{lema:C24-raiz-simultanea}
Para todo campo $k$ existe una extensión $k\subseteq K$ en la que cada polinomio no constante de $k[t]$ tiene al menos una raíz.
\end{lema}
\begin{proof}
Sea $\mathcal P$ el conjunto de polinomios mónicos no constantes de $k[t]$. Introducimos una indeterminada $t_f$ por cada $f\in\mathcal P$ y formamos
\[
\mathcal T=\set{t_f\given f\in\mathcal P},\qquad A=k[\mathcal T],\qquad
I=\sum_{f\in\mathcal P}(f(t_f))\subseteq A.
\]
El anillo de polinomios puede tener infinitas indeterminadas, pero cada polinomio usa sólo finitas y cada elemento del ideal es una suma finita de múltiplos de sus generadores.

El ideal $I$ es propio. De lo contrario existiría una identidad
\[
1=\sum_{i=1}^n a_i f_i(t_{f_i}),\qquad a_i\in A,
\]
con los $f_i$ distintos. Adjuntando sucesivamente una raíz de un factor irreducible de cada $f_i$, construimos una extensión $k\subseteq E$ con raíces $\alpha_i$ de todos ellos. Existe una evaluación $A\to E$ que envía $t_{f_i}$ a $\alpha_i$ y las demás indeterminadas a cero. Aplicada a la identidad da $1=0$, contradicción.

Por el resultado de existencia de ideales maximales de la clase 17, hay un maximal $\mathfrak m\supseteq I$. El cociente $K=A/\mathfrak m$ es un campo y el mapa de constantes $k\to K$ es inyectivo: su núcleo es un ideal propio del campo $k$. Para cada $f\in\mathcal P$, la clase de $t_f$ es una raíz, pues $f(t_f)\in\mathfrak m$. Dividir cualquier polinomio no constante por su coeficiente principal lo vuelve mónico sin cambiar sus raíces.
\end{proof}

La construcción proporciona raíces de los polinomios del campo de partida. En la extensión aparecen también polinomios con coeficientes nuevos: sobre $k=\mathbb Q$ podemos tener elementos $a,b$ con $a^2=2$ y $b^4=3$, y considerar ahora $t^2-a$. El lema por sí solo no afirma que todos los polinomios con esos nuevos coeficientes tengan raíz. Por eso iteramos el procedimiento:
\[
K_0=k\subseteq K_1\subseteq K_2\subseteq\cdots,
\qquad L=\bigcup_{n\geq0}K_n,
\]
donde cada polinomio no constante de $K_n[t]$ tiene una raíz en $K_{n+1}$.

\begin{prop}\label{prop:C24-union-cerrada}
La unión $L$ de esta cadena es un campo algebraicamente cerrado.
\end{prop}
\begin{proof}
Dos elementos de $L$ pertenecen simultáneamente a algún $K_n$. Su suma y producto se calculan allí, y el inverso de un elemento no nulo también pertenece a ese campo. Las inclusiones garantizan que las operaciones son compatibles, por lo que $L$ es un campo.

Si $f\in L[t]$ es no constante, sus finitos coeficientes pertenecen simultáneamente a algún $K_n$. Entonces $f$ tiene una raíz en $K_{n+1}\subseteq L$. Esto caracteriza a los campos algebraicamente cerrados.
\end{proof}

\begin{lema}\label{lema:C24-parte-algebraica}
Si $k\subseteq L$ y $L$ es algebraicamente cerrado, entonces
\[
\overline k=\set{\alpha\in L\given \alpha\text{ es algebraico sobre }k}
\]
es una clausura algebraica de $k$.
\end{lema}
\begin{proof}
El \cref{coro:C23-campo-algebraicos} muestra que $\overline k$ es un campo, algebraico sobre $k$ por su definición. Sea $\overline k\subseteq\overline k(\alpha)$ una extensión simple algebraica y sea $p$ el mínimo de $\alpha$. Como $L$ es algebraicamente cerrado, $p$ tiene una raíz $\beta\in L$. La adjunción de una raíz induce un homomorfismo
\[
\begin{tikzcd}[column sep=large,row sep=large]
\overline k\arrow[r,hook]\arrow[d,hook]&L\\
\overline k[t]/(p)\arrow[ur,hook,"\epsilon_\beta"',destacado]&
\end{tikzcd}
\]
que permite identificar $\overline k(\alpha)$ con $\overline k(\beta)\subseteq L$. La torre $k\subseteq\overline k\subseteq\overline k(\beta)$ es algebraica por el \cref{coro:C23-torre-algebraica}. Así $\beta$ es algebraico sobre $k$ y pertenece a $\overline k$, de donde la extensión simple inicial es trivial. Por el \cref{lema:C24-cerrado}, $\overline k$ es algebraicamente cerrado.
\end{proof}
El homomorfismo del diagrama queda determinado una vez elegida la raíz $\beta$. Sin fijar esa imagen puede haber varios homomorfismos que extiendan la inclusión del campo base. La unicidad de una clausura será salvo isomorfismo, sin exigir que el isomorfismo sea único.

%%%%%%%%%%%%%%%%%%%%%%%%%%%%%%%%%%%%%%%%%%%%%%%%%%%%%%%%%%%%%%%%%%%%%%%%%%%%%%%
\subsection{Prolongación de homomorfismos mediante Zorn}
%%%%%%%%%%%%%%%%%%%%%%%%%%%%%%%%%%%%%%%%%%%%%%%%%%%%%%%%%%%%%%%%%%%%%%%%%%%%%%%
% P05--P06.
\begin{lema}\label{lema:C24-inmersion}
Si $F/k$ es algebraica y $L$ es un campo algebraicamente cerrado que contiene a $k$, existe un homomorfismo $F\to L$ que restringe a la identidad de $k$.
\end{lema}
\begin{proof}
Consideremos el conjunto
\[
\mathcal Z=\set{(K,i_K)\given k\subseteq K\subseteq F,
\ i_K:K\to L,\ \restr{i_K}{k}=\id_k}.
\]
Es no vacío porque contiene la inclusión de $k$. Lo ordenamos por prolongación: $(K,i_K)\preceq(K',i_{K'})$ si $K\subseteq K'$ y $\restr{i_{K'}}{K}=i_K$.

Para una cadena no vacía, sea $K_C$ la unión de sus campos. Es un campo y las aplicaciones se reúnen en un homomorfismo $i_C:K_C\to L$: el valor de un elemento se calcula en cualquier miembro que lo contenga, y la compatibilidad asegura que es independiente de la elección. Para comprobar suma y producto basta tomar un miembro que contenga los dos elementos. Esta aplicación es una cota superior de la cadena. La cadena vacía admite como cota la inclusión de $k$. El esquema de compatibilidad es
\[
\begin{tikzcd}[column sep=large,row sep=large]
&L&\\
K\arrow[r,hook]\arrow[ur,"i_K"]&
K'\arrow[r,hook]\arrow[u,"i_{K'}"']&
K_C\arrow[ul,dashed,"i_C"',destacado].
\end{tikzcd}
\]
Por Zorn existe un elemento maximal $(G,i_G)$. Sea $H=i_G(G)$; el mapa $i_G:G\to H$ es un isomorfismo porque todo homomorfismo de campos es inyectivo.

Si $G\ne F$, elijamos $\alpha\in F\setminus G$. Es algebraico sobre $G$, pues ya lo es sobre $k$. Su mínimo $g\in G[t]$ se transforma mediante el isomorfismo de coeficientes en un irreducible $h\in H[t]$. El campo $L$ contiene una raíz $\beta$ de $h$. Por el transporte de extensiones simples de la clase 22, $i_G$ se prolonga a un isomorfismo $G(\alpha)\to H(\beta)\subseteq L$ que envía $\alpha$ a $\beta$:
\[
\begin{tikzcd}[column sep=large,row sep=large]
G\arrow[r,"i_G","\sim"']\arrow[d,hook]&
H\arrow[d,hook]\\
G(\alpha)\arrow[r,dashed,"\widetilde i_G","\sim"',destacado]&
H(\beta)\arrow[r,hook]&L.
\end{tikzcd}
\]
Como $G(\alpha)\subseteq F$ contiene propiamente a $G$, contradice la maximalidad. Luego $G=F$.
\end{proof}

\begin{teo}\label{teo:C24-clausura}
Todo campo admite una clausura algebraica. Dos clausuras de $k$ son isomorfas mediante un isomorfismo que fija a $k$.
\end{teo}
\begin{proof}
Para la existencia, construimos un campo algebraicamente cerrado $L\supseteq k$ mediante la cadena anterior y tomamos su parte algebraica sobre $k$, que es una clausura por el \cref{lema:C24-parte-algebraica}.

Si $\overline k_1$ y $\overline k_2$ son dos clausuras, el \cref{lema:C24-inmersion} proporciona un homomorfismo $i:\overline k_1\to\overline k_2$ que fija a $k$. Su imagen $H$ es algebraicamente cerrada, por ser isomorfa a $\overline k_1$. La extensión $\overline k_2/H$ es algebraica: cada elemento ya es algebraico sobre $k\subseteq H$. El \cref{lema:C24-cerrado} obliga a que $H=\overline k_2$, de modo que $i$ es sobreyectivo y es un isomorfismo.
\end{proof}
En particular, $\overline{\mathbb Q}$ es una clausura algebraica del campo racional.

%%%%%%%%%%%%%%%%%%%%%%%%%%%%%%%%%%%%%%%%%%%%%%%%%%%%%%%%%%%%%%%%%%%%%%%%%%%%%%%
\section{Campos de descomposición y normalidad}
%%%%%%%%%%%%%%%%%%%%%%%%%%%%%%%%%%%%%%%%%%%%%%%%%%%%%%%%%%%%%%%%%%%%%%%%%%%%%%%
%%%%%%%%%%%%%%%%%%%%%%%%%%%%%%%%%%%%%%%%%%%%%%%%%%%%%%%%%%%%%%%%%%%%%%%%%%%%%%%
\subsection{Existencia, unicidad y cota del grado}
%%%%%%%%%%%%%%%%%%%%%%%%%%%%%%%%%%%%%%%%%%%%%%%%%%%%%%%%%%%%%%%%%%%%%%%%%%%%%%%
% P07--P08.
Una familia de polinomios se \emph{descompone} sobre una extensión cuando cada uno se factoriza en términos lineales. Para una familia finita de polinomios no nulos, esto equivale a que se descomponga su producto. Por eso podemos concentrarnos en un solo polinomio.

\begin{defi}
Sea $0\ne f\in k[t]$ de grado $d$. Un \emph{campo de descomposición} de $f$ sobre $k$ es una extensión $F/k$ tal que
\[
f(t)=c\prod_{i=1}^d(t-\alpha_i)\quad\text{en }F[t],
\qquad F=k(\alpha_1,\ldots,\alpha_d),
\]
con $c\in k^\times$. Las raíces se cuentan con sus multiplicidades; para $d=0$ se toma $F=k$.
\end{defi}

\begin{prop}\label{prop:C24-descomposicion}
Todo polinomio no nulo $f\in k[t]$ admite un campo de descomposición $F$, único salvo isomorfismo sobre $k$, y $[F:k]\leq(\deg f)!$.

Más generalmente, si $\iota:k'\to k$ es un isomorfismo y $g\in k'[t]$ se transforma en $f$ al aplicar $\iota$ a sus coeficientes, $\iota$ se extiende a un isomorfismo entre cualquier campo de descomposición de $g$ y cualquier campo de descomposición de $f$.
\end{prop}
\begin{proof}
Probamos existencia y cota por inducción sobre $d=\deg f$. Para $d=0$ o $d=1$, basta $F=k$ y $[F:k]=1$. Si $d>1$, tomamos un factor irreducible $q$ de $f$. El campo $E=k[t]/(q)=k(\alpha)$ contiene una raíz de $f$ y satisface $[E:k]=\deg q\leq d$. En $E[t]$ escribimos $f=(t-\alpha)h$, con $\deg h=d-1$. Por inducción, $h$ tiene un campo de descomposición $F/E$ de grado a lo sumo $(d-1)!$. El campo $F$ está generado sobre $k$ por $\alpha$ y las raíces de $h$, luego es un campo de descomposición de $f$. Además,
\[
[F:k]=[F:E][E:k]\leq(d-1)!\,d=d!.
\]

Para la unicidad general, fijamos una clausura $\overline k$. Un campo de descomposición $F/k$ se puede incluir en ella por el \cref{lema:C24-inmersion}, pues es algebraico. Si $G/k'$ es un campo de descomposición de $g$, el mismo lema, identificando $k'$ con $k$ mediante $\iota$, da una inclusión $j:G\to\overline k$ que extiende a $\iota$. Sus imágenes de las raíces de $g$ son todas las raíces de $f$, porque transforman una factorización lineal completa de grado $d$ en otra del mismo grado. Por tanto $j(G)$ es el subcampo de $\overline k$ generado por $k$ y las raíces de $f$. Ése es también $F$, así que $j:G\to F$ es un isomorfismo. Los dos transportes se resumen en
\[
\begin{tikzcd}[column sep=large,row sep=large]
k'[t]\arrow[r,"\widehat\iota","\sim"']&k[t]\\
k'\arrow[r,"\iota","\sim"']\arrow[u,hook]&k\arrow[u,hook]
\end{tikzcd}
\qquad
\begin{tikzcd}[column sep=large,row sep=large]
G\arrow[r,dashed,"j","\sim"',destacado]&F\\
k'\arrow[r,"\iota","\sim"']\arrow[u,hook]&k\arrow[u,hook].
\end{tikzcd}
\]
El caso $k'=k$ e $\iota=\id$ prueba la unicidad sobre el campo base.
\end{proof}

%%%%%%%%%%%%%%%%%%%%%%%%%%%%%%%%%%%%%%%%%%%%%%%%%%%%%%%%%%%%%%%%%%%%%%%%%%%%%%%
\subsection{Ejemplos y extensiones normales}
%%%%%%%%%%%%%%%%%%%%%%%%%%%%%%%%%%%%%%%%%%%%%%%%%%%%%%%%%%%%%%%%%%%%%%%%%%%%%%%
% P09.
\begin{ejem}
El campo de descomposición de $t^2+1$ sobre $\mathbb Q$ es $\mathbb Q(i)$, mientras que sobre $\mathbb R$ es $\mathbb C$. En ambos casos las raíces son $i$ y $-i$, pero se generan campos sobre bases distintas.
\end{ejem}

\begin{ejem}[Raíces octavas de la unidad]\label{ejem:C24-octavas}
Sea $\zeta=e^{2\pi i/8}$. Las raíces de $t^8-1$ son $1,\zeta,\ldots,\zeta^7$, de modo que su campo de descomposición sobre $\mathbb Q$ es $\mathbb Q(\zeta)$. En la factorización
\[
t^8-1=(t^4-1)(t^4+1),
\]
el primer factor tiene raíces $1,-1,i,-i$ y el segundo tiene las potencias impares de $\zeta$, que se destacan en la figura.
\begin{center}
\begin{tikzpicture}[scale=1.65,>=Stealth]
\draw[gray] (0,0) circle (1);
\draw[gray] (-1.2,0)--(1.65,0);
\draw (0,0)--(45:1);
\draw[dashed] (45:1)--({sqrt(2)},0)--(-45:1);
\draw[dashed] (0,0)--(-45:1);
\foreach \j in {0,...,7}{\fill ({45*\j}:1) circle (1.1pt);}
\foreach \j in {1,3,5,7}{\draw[rojoDiagramaCIMAT,line width=.8pt] ({45*\j}:1) circle (.09);}
\fill (0,0) circle (1.1pt) node[below left=2pt] {$0$};
\fill ({sqrt(2)},0) circle (1.1pt) node[below=5pt] {$\sqrt2$};
\node[above right=4pt] at (45:1) {$\zeta$};
\node[above left=4pt] at (1,0) {$1$};
\node[left=4pt] at (-1,0) {$-1$};
\node[above=4pt] at (0,1) {$i$};
\node[below=4pt] at (0,-1) {$-i$};
\end{tikzpicture}
\end{center}

El polinomio $t^4+1$ es irreducible sobre $\mathbb Q$: tras el cambio $t\mapsto t+1$ resulta
\[
(t+1)^4+1=t^4+4t^3+6t^2+4t+2,
\]
al que se aplica Eisenstein con el primo $2$. El cambio de variable es un automorfismo de $\mathbb Q[t]$, por lo que conserva irreducibilidad. Así $t^4+1$ es el mínimo de $\zeta$, y $\mathbb Q(\zeta)$ es ya su campo de descomposición, con
\[
[\mathbb Q(\zeta):\mathbb Q]=4.
\]
El grado cuatro mejora ampliamente las cotas $8!$ y $4!$ de la proposición.
\end{ejem}

\begin{defi}
Una extensión $F/k$ es \emph{normal} si todo irreducible de $k[t]$ que tiene una raíz en $F$ se descompone por completo sobre $F$.
\end{defi}

\begin{teo}\label{teo:C24-normal}
Una extensión $F/k$ es finita y normal si y sólo si es un campo de descomposición de algún polinomio no nulo de $k[t]$.
\end{teo}
\begin{proof}
Si la extensión es finita, cualquier base vectorial es también una familia de generadores como campo. Escribamos $F=k(\alpha_1,\ldots,\alpha_r)$ y sean $p_i\in k[t]$ sus polinomios mínimos. Por normalidad, todos los $p_i$ se descomponen en $F$. El producto $f=\prod_i p_i$ se descompone y sus raíces generan $F$: pertenecen a $F$ y entre ellas están todos los $\alpha_i$.

Recíprocamente, sea $F$ un campo de descomposición de $f$. Es finito por la \cref{prop:C24-descomposicion}. Incluyámoslo en $\overline k$ y sea $p\in k[t]$ irreducible con una raíz $\alpha\in F$. Si $\beta\in\overline k$ es cualquier otra raíz de $p$, existe un isomorfismo $\iota:k(\alpha)\to k(\beta)$ que fija a $k$ y envía $\alpha$ a $\beta$. El campo $F$ es de descomposición de $f$ sobre $k(\alpha)$, y $F(\beta)$ lo es sobre $k(\beta)$. La unicidad general de la proposición prolonga $\iota$ a un isomorfismo $F\to F(\beta)$ que fija a $k$. Por tanto ambos tienen la misma dimensión finita sobre $k$. Como además $F\subseteq F(\beta)$, coinciden y $\beta\in F$. Todas las raíces de $p$ pertenecen a $F$, lo que prueba la normalidad.
\end{proof}

\begin{obs}[Referencias de lectura]
Para clausuras algebraicas, véase \citet[cap.~VII, apartado 2.1, lema 2.1, definición 2.2, teorema 2.3 y resultados 2.4--6/8]{aluffi2009}; la comprobación sobre sus raíces corresponde al ejercicio VII.2.2. Los campos de descomposición y la normalidad se tratan en VII.4.1--4/7--8. La clase 25 retoma la caracterización de normalidad y desarrolla su diagrama de comparación de raíces.
\end{obs}
